\section{Valutazione del Design}

\subsection{Valutazione Euristica e Cognitive Walkthrough}
% Descrivere la tecnica di valutazione del design adottata per valutare il prototipo Figma (es. Valutazione Euristica o Cognitive Walkthrough), immedesimandosi nelle personas individuate negli assignment precedenti.

\subsection{Risultati della Valutazione e Confronto}
% Riassumere i risultati emersi dalla valutazione ed effettuare un confronto.

\subsection{Lista delle Modifiche Prioritarie}
% Inserire la lista ordinata e con relative priorità degli interventi da fare per migliorare il prototipo prima di procedere all'implementazione.

\begin{table}[H]
    \centering
    \small
    \renewcommand{\arraystretch}{1.3}
    \begin{tabular}{
        p{1.5cm}
        |>{\raggedright\arraybackslash}p{6cm}
        |>{\centering\arraybackslash}p{2.5cm}
        |>{\centering\arraybackslash}p{2.5cm}
    }
    \toprule
    \rowcolor{culturiaRed}
    \textcolor{culturiaWhite}{\textbf{ID}} &
    \textcolor{culturiaWhite}{\textbf{Descrizione Intervento}} &
    \textcolor{culturiaWhite}{\textbf{Gravità}} &
    \textcolor{culturiaWhite}{\textbf{Priorità}} \\
    \midrule
    \rowcolor{culturiaBeige}
    \textbf{INT-01} & & & \\
    
    \textbf{INT-02} & & & \\
    
    \rowcolor{culturiaBeige}
    \textbf{INT-03} & & & \\
    \bottomrule
    \end{tabular}
    \caption{Lista delle modifiche da effettuare prima dell'implementazione.}
    \label{tab:priorita_modifiche}
\end{table}
